% !TEX root = ../thesis.tex

% soft colors (print-friendly)
\definecolor{baseline}{RGB}{255,230,200} % light orange
\definecolor{newmodel}{RGB}{220,235,255} % light blue
% define colorbox for the baseline e new model



\label{ch:experiments}
\section{Introduction}
This chapter describes the hyperparameter optimization procedure employed to fairly compare different reservoir computing architectures across multiple benchmark tasks. 
%We utilized Bayesian optimization via the Optuna framework to systematically explore the hyperparameter space for each model
%variant while maintaining a fixed parameter budget of approximately 100,000 reservoir parameters.
In our empirical evaluation process aimed to assess the performances of the architectures proposed in Chapter \ref{ch:methodology}.
We consider (i) memory based tasks such as: short-term memory capacity (MC), long temporal dependacies with sequential classification bechmarks (Adiac, FordA, FordB, sMNIST, psMNIST and npCIFAR-10),
 and (ii) supervised benchmarks spanning synthetic temporal regression (NARMA, Mackey--Glass, Lorenz96).
\section{Memory Tasks}
Memory is a fundamental aspect of recurrent neural networks, the capability of retaining and utilizing past information for currently processing inputs
is crucial for tasks involving temporal dependencies.
\subsection{Memory Capacity}
\label{sec:exp:mc}
Memory Capacity (MC) taks was introduced by Jeager in \cite{memorycapacity}, to quantify the capacity of recurrent networks to represent past inputs by correlating the past events of an i.i.d. input
stream with the network output. We assess short-term memory as follows:

Let $u(t)\sim \mathcal{U}[-0.8,0.8]$ be an i.i.d. input stream driving the system of discrete time steps $T=6000$ and let $x(t)\in\mathbb{R}^N$ be the state at time $t$.
For each delay $k\in\{1,\dots,K\}$ we train a linear readout $\hat{u}(t-k) = w_k^\top x(t)$ 
to reconstruct the input $u(t-k)$ from the current state $x(t)$.
 We define the memory capacity at delay $k$ as the squared correlation coefficient between the target $u(t-k)$ and the reconstructed signal $\hat{u}(t-k)$:

\begin{equation}
\mathrm{MC}_k \;=\; \frac{\mathrm{cov}^2\!\left(u(t-k),\hat{u}(t-k)\right)}{\mathrm{var}\!\left(u(t-k)\right)\,\mathrm{var}\!\left(\hat{u}(t-k)\right)} \in [0,1],
\end{equation}
and the total memory capacity as
\begin{equation}
\mathrm{MC} \;=\; \sum_{k=1}^{\infty} \mathrm{MC}_k.
\label{eq:mc_k}
\end{equation}

\paragraph{Experimental Setting}
In order to compute and compare MC of our models with others, we adopted the same setting widely used in RC community, using a reservoir size of $N_{r}=100$ units linearly activated and trained by Ridge Regression, 
setting the total lag $k=200$ (twice reservoir size) from equation \ref{eq:mc_k}, the reason behind this choice is that the maximum
MC achievable by an N-unit linear reservoir is bounded by $N$, a proof is provided by \cite[Proposition 2.]{memorycapacity}.

For the warmup period, we discarded the first $T_{\mathrm{wash}}=100$ timesteps to remove dependacies on initial conditions.  
The readout weights are trained on $5000$ timesteps and then we evaluted the test MC on the last $900$ timesteps.
Since we are interested to exactly recall the target signal, we set the regularization parameter $\lambda=0$ in the Ridge Regression.

For the remaining hyperparameters, we performed a random search to estimate the most important parameters affecting the MC performance, often
being spectral radious and input scaling. We left the leakage rate $\alpha = 1.0$.

Our experiment were repeated over different number of layer by keeping the same number of total reservoir neurons, so 1 layer with $N_{r}=100$ units, 2 layers with $N_{r}=50$ units each, 4 layers with $N_{r}=25$ and so on.


\paragraph{Reported results}
We report $\mathrm{MC}$ and optionally the curve $\mathrm{MC}_k$ vs.\ $k$ to characterize how memory decays with delay.

\subsection{Results}
In order to asses if adding modularity in deep reservoir architecture can improve performances in terms of memory capacity,
we compare our results between the best configuration for each architecture.

%add table here % with ARCHITECTURE, MEMORY CAPACITY TRAIN/TEST
\begin{table}[htbp!]
\centering
\caption{Memory Capacity results for different architectures averaged over 3 trials from the best performing model. The higher the better.}
\label{tab:mc_results}
\begin{tabular}{@{}lcc@{}}
\toprule
\textbf{Architecture} & \textbf{Train MC} & \textbf{Test MC} \\
\midrule
DeepESN & 55.8 $\pm$ 2.64 & 53.01 $\pm$ 2.45  \\
DCR & 99.1 $\pm$ 0.1 & \textbf{99} $\pm$ 0.1 \\
DAR & $54.65 \pm 2.7$  & $51.83 \pm 2.8$  \\  
% add line here for different architectures
\midrule 
DeepRON & $49.14 \pm 2.2$  & $43.25 \pm 3.1$   \\    
DCO & 41.9 $\pm$ 2.13 & 38.86 $\pm$ 2.52 \\
DAO & 48.16 $\pm$ 2.16 & 44.5 $\pm$ 3.8 \\
\bottomrule
\end{tabular}
\end{table}

The results in Table \ref{tab:mc_results} show that with the proposed DCR we are able to achieve the same performance of the SCR, this is true because in the case of linear DCR with 1 unit per layer we get exactly
the same architecture of SCR.
The other variant with antisymmetric coupling, instead, doesn't achieve very high performances in terms of memory MC,
 with respect to the DeepESN. 
 For the Oscillator based reservoirs, we can't recover the same behaviour in case of cycle architecture as for DeepESN, this is probably due to the fact that the oscillators dynamics need to be carefully tuned layerwise
 to achieve higher performances. 
 In DeepRONs we can control the eigenvalues location by tuning the $\gamma$ and $\epsilon$ parameters, but the system becomes good only in a small regime of parameters, so further investigation is needed
 on a eterogeneous setting of deep oscillator reservoirs. 

In the figure \ref{fig:mc_dcr_layers} below we can see the MC curves for DCR as we increase the number of layers, in this example we fixed $\rho=0.5$ to ensure contractivity, for what we said in \ref{sec:method:cycle_architecture} a too high 
spectral radius can lead to a global system with much higher spectral radius. \todo{metto un immagine anche per vedere cosa succede quando aumento man mano il raggio spettrale finche non va oltre 1?}
The 100 units in cycle perfectly recall the input up to delay 99, as soon we halve the number of layers having 2 units in each layer (yellow line in the figure), the performances drop but still remain high, yielding a test of $MC = 67.3$.

\begin{figure}[htbp!]
    \centering
    \includegraphics[width=0.9\textwidth]{images/memory_capacity_vs_delay_esn.png}
    \caption{Memory Capacity curves for DCR with different number of layers.}
    \label{fig:mc_dcr_layers}
\end{figure}

Finally considering the case where we take in "edge of chaos" regime, the DCR architecture, by setting a spectral radius of about $\rho = 0.7$, we get a $\rho_\text{global} \approx 0.99$, in the figure below we 
can see both MC curve and a plot of the Jacobian eigenvalues distribution.
% insert figure with the curve up and below the eigenvalues give a lot of space to both the images
\begin{figure}[htbp!]
    \centering
    \includegraphics[width=0.9\textwidth]{images/memory_capacity_vs_delay_DCR_Edgeofchaos.png}
    \caption{Memory Capacity curve for DCR in edge of chaos regime.}
    \label{fig:mc_dcr_eoc}
\end{figure}

% eigenvalues
\begin{figure}[htbp!]
    \centering
    \includegraphics[width=0.9\textwidth]{images/eigenvalue_distributions_esn_edgeofchaos.png}
    \caption{Jacobian eigenvalues distribution for DCR in edge of chaos regime.}
    \label{fig:mc_dcr_eoc_eigenvalues}
\end{figure}

\section{Long temporal dependencies tasks}  
This set of experiments aims to evaluate the ability of 
different reservoir to capture in different ways long temporal dependancies,
 we consider three different sequential image classification 
tasks on two dataset MNIST and CIFAR-10.


\begin{comment}
We cast image classification as a sequence modeling problem by feeding pixels (or small patches) sequentially.

\subsection{MNIST}
\label{sec:exp:mnist}

MNIST images are $28\times 28$ grayscale. We consider:
\begin{itemize}
    \item \textbf{sMNIST (sequential MNIST):} pixels are streamed row-wise as a length-$784$ sequence, $u(t)\in[0,1]$.
    \item \textbf{pMNIST (permuted sequential MNIST):} a fixed random permutation $\pi$ of $\{1,\dots,784\}$ is applied to the pixel order, yielding $u(t)=\mathrm{vec}(I)_{\pi(t)}$. The permutation is fixed across all samples and splits.
\end{itemize}
The label is provided at the end of the sequence; the model output is taken from the final state (or pooled states) and passed through a classifier.

\paragraph{Metric.}
We report test accuracy (\%).

\subsection{npCIFAR-10 (sequential CIFAR-10 variant)}
\label{sec:exp:npcifar}

CIFAR-10 images are $32\times 32$ with 3 channels. We define a sequential variant tailored to stress long-range dependencies:
\begin{itemize}
    \item \textbf{Sequence construction:} flatten the image into a sequence of length $32\cdot 32=1024$ tokens, where each token is a 3D vector (RGB) or a scalar after luminance projection.
    \item \textbf{Permutation (P):} apply a fixed random permutation of token indices to destroy locality.
    \item \textbf{Noise (N):} optionally add i.i.d.\ Gaussian noise $\epsilon\sim\mathcal{N}(0,\sigma^2)$ to inputs during training (and optionally testing) to assess robustness.
\end{itemize}
We refer to this setting as \emph{npCIFAR-10}. Exact choices (RGB vs.\ scalar, permutation seed, noise level $\sigma$) are specified in Sec.~\ref{sec:exp:details}.
\end{comment}

\subsection{Sequential MNIST (sMNIST)}
The Sequential MNIST task treats each $28 \times 28$ grayscale image as a sequence of 784 scalar values, by flattening the image matrix,
 thus destorying spacial information in favor of temporal structure. The pixels are presented in a fixed order (row-wise). Data is presented one pixel at a time, so the model must classify the digit (0--9) after observing the entire sequence. 
 This task evaluates a model's ability to retain information over long temporal horizons while processing one-dimensional sequential inputs.
% add example of smnist image as sequence

\begin{figure}[ht!]
    \centering
    \includegraphics[width=0.7\textwidth]{images/smnist.png}
    \caption{Example of Sequential MNIST (sMNIST) input sequence. Each pixel of the $28 \times 28$ image is presented sequentially as a scalar input over 784 timesteps.}
    \label{fig:smnist_sequence}
\end{figure}

\subsection{Permuted Sequential MNIST (psMNIST)}
The Permuted Sequential MNIST task applies a fixed random permutation to the pixel ordering of sMNIST sequences. This destroys the spatial structure of the original images, making the task purely sequential and significantly more challenging. The permutation is generated deterministically using a fixed random seed to ensure reproducibility.

\subsection{Noise-Padded CIFAR-10 (npCIFAR-10)}
The Noise-Padded CIFAR-10 task consists to take CIFAR-10 dataset that contains $32 \times 32 \times 3$ color images across 10 classes. 
Whe first two dimensions representing spatial dimensions and the third dimension representing RGB color channels.
Each image is flattened row-wise into a sequence representation of 32 rows, where each row contains 96 values (32 pixels $\times$ 3 channels).
We fed 96-dimensional vectors sequentially over 32 timesteps, then for the remaining 968 timesteps we add uniform random noise in the range $[-1, 1]$, resulting in sequences of length 1000. 
This setup tests the model's ability to extract and maintain relevant information while being exposed to subsequent irrelevant noise.

\subsection{Results}
In the table below we report results for the three sequential image classification tasks, comparing the best test accuracy over 3 trials achieved by each architecture across the datasets.
% use ligh color blue for evidencing new model put an underline for the best result in each column
% use light color orange for evidencing the baselines

\begin{table}[ht!]
    \centering
    \caption{Performance Summary. Comparison of the best test accuracy over 3 trials achieved by each architecture across the three datasets.
    We higlight in \ \colorbox{newmodel}{\strut proposed modular models}\ and\ \colorbox{baseline}{\strut shallow and deep models}.
    With $\pm$ we indicate the standard deviation across trials.}
    \label{tab:performance_summary}
    \begin{tabular}{lccc}
        \toprule
        \textbf{Architecture} & \textbf{sMNIST} & \textbf{psMNIST} & \textbf{npCIFAR-10} \\
        \midrule
        \rowcolor{baseline}
        ESN                 &$82.93\% \pm 0.89$& $78.87\% \pm 1.82$ & $32.06\% \pm 0.19$ \\
        \rowcolor{baseline}
        DeepESN             &$85.94\% \pm 0.46$& $ 80.62\% \pm 0.95$ & $\underline{\mathbf{39.83\%}} \pm 0.19$ \\
        \rowcolor{newmodel}
        DCR     & $86.81\% \pm 0.49$ & $75.92\% \pm 0.90$ & $32.28\% \pm 0.19$ \\
        \rowcolor{newmodel}
        DCR (No Recurrence) & $86.32\% \pm 0.35$  & $79.31\% \pm 1.62$ & $33.91\%  \pm 0.39$ \\
        \rowcolor{newmodel}
        DAR & $\underline{\textbf{88.29\%}} \pm 0.89$ & $\underline{\textbf{81.33\%}} \pm 0.65$ & $34.37\% \pm 0.16$ \\
        \midrule 
        \midrule
        \rowcolor{baseline}
        RON                 & $\underline{\textbf{90.7\%}} \pm 0.51$ & $\underline{\textbf{86.5\%}} \pm 1.23$ & $28.42\% \pm 0.21$ \\
        \rowcolor{baseline}
        DeepRON             & $84.8\% \pm 1.4$ & $82.1\% \pm 4.5$ & $26.82\% \pm 0.1$ \\
        \rowcolor{newmodel}
        DCO     & $89.17\% \pm 0.8$ & $51.03\% \pm 0.8$ & $26.36\% \pm 0.23$ \\
        \rowcolor{newmodel}
        DCO (No Recurrence) & $61.25\% \pm 0.51$ & $39.6\% \pm 1.4$ & $26.87\% \pm 0.12$ \\
        \rowcolor{newmodel}
        DAO & $88.7\% \pm 0.62$ & $83.32\% \pm 2.1$ & $\underline{\textbf{32.12\%}} \pm 0.15$ \\
        \bottomrule
    \end{tabular}
\end{table}

\todo{discutere risultati}
- DeepESN seems to be better thna other architectures on npCIFAR-10, this can be attributed to the
 fact that in DeepESNs layer act as filter banks, so they can ignore all the irrelevant frequency components 
 introduced by noise padding, focussing only on the relevant part of the sequence that contains the image information.
 The reason behind the worse performance of modules architectures on the task could be also attributed in how the information is propagated,
 probably the inter-layer connections trade the irrelevant information across timestep (the noise padding part), making more difficult 
 to extract the relevant features at the first steps of the sequence.
 - However on sMNIST ans psMNIST we can se that the antisymmetric coupling improves performances, also the DCR with and without recurrence 
 show improvements with respect to DeepESN on the first two tasks

\section{Time Series Classification and Regression Tasks}

\subsection{FordA and FordB}
The Ford datasets are binary classification tasks from the UCR Time Series Archive. FordA contains 3601 training and 1320 test samples, while FordB contains 3636 training and 810 test samples. Each sample is a univariate time series of length 500, representing sensor measurements from automotive engine systems. The task is to classify whether a specific symptom is present.

\subsection{Mackey-Glass Time Series Prediction}
The Mackey-Glass (MG) system is a classical chaotic time series benchmark defined by the delay differential equation:
\begin{equation}
\frac{dx(t)}{dt} = \beta\frac{x(t-\tau)}{1+x(t-\tau)^n} - \gamma x(t),
\label{eq:mg}
\end{equation}
We use standard parameters that produce chaotic dynamics ($\tau = 17$, $n = 10$, $\beta = 0.2$, $\gamma = 0.1$). The task is 84-step ahead prediction, i.e., given $x(t)$, predict $x(t+84)$.
We employ a washout period of 200 timesteps to allow the reservoir state to stabilize. Performance is measured using the Normalized Root Mean Squared Error (NRMSE):
\begin{equation}
\text{NRMSE} = \frac{\sqrt{\frac{1}{N}\sum_{i=1}^{N}(\hat{y}_i - y_i)^2}}{\sigma_y}
\end{equation}
where $\sigma_y$ is the standard deviation of the target signal.

\subsection{Results}


% make a table for fordA and fordb
\begin{table}[ht!]
    \centering
    \caption{FordA and FordB classification accuracy (\%) for different architectures.}
    \label{tab:ford_results}
    \begin{tabular}{lcc}
        \toprule
        \textbf{Architecture} & \textbf{FordA Accuracy} & \textbf{FordB Accuracy} \\
        \midrule
        ESN                 & $ 67\% \pm 0.2$ & $57.32\% \pm 0.146\%$ \\
        DeepESN             & $75.2\% \pm 0.016$ & $0.65 \pm 0.13$ \\
        DCR     & $0 \pm 0.9$ & $0 \pm 1.1$ \\
        DCR (No Recurrence) & $4$ & $0.4 \pm 1.2$ \\
        DAR & $\mathbf{0.6929 \pm 0.02}$ & $\mathbf{0 \pm 1.0}$ \\
        \midrule
        RON                 & $0.69 \pm 0.02 $ & $0.73 \pm 0.04$ \\
        DeepRON             & $66.3\% \pm 0.11 $ & $0.5 \pm 1.3$ \\
        DCO     & $65.2\% \pm 0.3$ & $0. \pm 1.1$ \\
        DCO (No Recurrence) & $0.55 \pm 0.03$ & $0.4 \pm 1.2$ \\
        DAO & $\mathbf{0.6821 \pm 0.02}$ & $\mathbf{0.5 \pm 1.0}$ \\
        \bottomrule
    \end{tabular}
\end{table}

Commenting results about FordA and FordB...
Lorem ipsum dolor sit amet, consectetur adipiscing elit, sed do eiusmod tempor
 incididunt ut labore et dolore magna aliqua. Ut enim ad minim veniam, 
 quis nostrud exercitation ullamco laboris nisi ut aliquip ex ea commodo consequat. 
 Duis aute irure dolor in reprehenderit in voluptate velit esse cillum dolore eu fugiat nulla pariatur.
  Excepteur sint occaecat cupidatat non proident...
% make a table for mackey glass
\begin{table}[ht!]
    \centering
    \caption{MG84 prediction NRMSE for different architectures. Lower is better. }
    \label{tab:mackeyglass_results}
    \begin{tabular}{lc}
        \toprule
        \textbf{Architecture} & \textbf{MG84 NRMSE} \\
        \midrule
        ESN                 & $3\times 10^{-2} \pm 2.4\times 10^{-3}$ \\
        DeepESN             & $0.198 \pm 0.009$ \\
        DCR     & $0.205 \pm 0.01$ \\
        DCR (No Recurrence) & $0.210 \pm 0.01$ \\
        DAR & $\mathbf{0.185 \pm 0.008}$ \\
        \midrule
        RON                 & $1.8 \times 10^{-2} \pm 6\times 10^{-3}$\\
        DeepRON             & $0.215 \pm 0.011$ \\
        DCO     & $0.220 \pm 0.01$ \\
        DCO (No Recurrence) & $0.218 \pm 0.01$ \\
        DAO & $\mathbf{0.1850 \pm 0.009}$ \\
        \bottomrule
    \end{tabular}
\end{table}
\clearpage
\section{Experimental details}

\subsection{Bayesian Optimization Framework}

For the rest of the experiments we employed Optuna \cite{optuna_2019} for Bayesian hyperparameter optimization.
 Optuna uses the Tree-structured Parzen Estimator (TPE) \cite{bergstra2011algorithms}
  as its default sampling strategy, which we utilized for the heavy hyperparameter search involved in our experiments, given the great number of parameters to tune for each architecture.

The TPE algorithm models the conditional probability $p(x|y)$ of hyperparameters $x$ given the objective value $y$ using two density estimators:
\begin{equation}
p(x|y) = \begin{cases}
\ell(x) & \text{if } y < y^* \\
g(x) & \text{if } y \geq y^*
\end{cases}
\end{equation}
where $y^*$ is a quantile threshold (typically the best 25\% of observations), $\ell(x)$ 
is the density formed by observations with objective values below $y^*$, and $g(x)$ is formed by the remaining observations.
 New hyperparameter configurations are sampled by maximizing the ratio $\ell(x)/g(x)$,
  which approximates the Expected Improvement acquisition function.

\subsection{Hyperparameter Search Spaces}

\subsubsection{Echo State Network Hyperparameters}

Table \ref{tab:esn_hyperparams} presents the hyperparameter search space for ESN-based models: DeepESN, DCR and DAR.

\begin{table}[ht!]
\centering
\caption{All parameters with log scale are sampled from a log-uniform distribution\footnote{We use a log-uniform distribution to ensure uniform sampling in the logarithmic space, which allows for a more balanced exploration of hyperparameters across several orders of magnitude.}.}
\label{tab:esn_hyperparams}
\begin{tabular}{@{}lccc@{}}
\toprule
\textbf{Hyperparameter} & \textbf{Range} & \textbf{Scale} & \textbf{Description} \\
\midrule
$\rho$ & $[0.1, 9.0]$ & Log &  Spectral radius \\
$\omega_{\text{in}}$ & $[0.1, 1.0]$ & Log &  Input scaling \\
$\alpha$ & $[0.001, 1.0]$ & Log & Leaky integrator coefficient \\
$\gamma$ & $[0.0001, 0.5]$ & Log & Diffusion regularizer \\
$\epsilon$ & $20.0$ & Fixed & Inter-layer coupling strength \\
Number of layers & $\{1\}$ or $\{5, 10\}$ & Categorical & Architecture-dependent \\
$\omega_{\text{inter}}$ & $[0.1, 1.0]$ & Fixed & Inter-layer scaling \\
\bottomrule
\end{tabular}
\end{table}

\subsubsection{Randomized Oscillators Network Hyperparameters}

Table \ref{tab:ron_hyperparams} presents the hyperparameter search space for RON-based models: DeepRON, DCO and DAO.

\todo{questa tabella forse posso togliere la descrizione e metterla anche nell'appendice...}
\begin{table}[ht!]
\centering
\caption{RON hyperparameter search space.}
\label{tab:ron_hyperparams}
\begin{tabular}{@{}lccc@{}}
\toprule
\textbf{Hyperparameter} & \textbf{Range} & \textbf{Scale} & \textbf{Description} \\
\midrule
$\tau$ & $[10^{-4}, 1.0]$ & Log & Integration time step \\
$\rho$ & $[0.1, 9.0]$ & Log & Spectral Radius \\
$\omega_{\text{in}}$ & $[0.05, 2.0]$ & Log & Input scaling \\
$\epsilon_c$ & $[0.01, 5.0]$ & Log & Inter-layer coupling \\
$\gamma_c$ & $[0.0001, 0.5]$ & Log & Diffusive strength \\
$\gamma$ (damping) & $1.0 \pm 0.05$ & Fixed &  Damping  \\
$\epsilon$ (stiffness) & $1.0 \pm 0.05$ & Fixed &  Stiffness \\
Number of layers & $\{5, 10\}$ & Categorical & Architecture-dependent \\
$\omega_{\text{inter}}$ & $[0.1, 1.0]$ & Fixed & Inter-layer scaling \\
\bottomrule
\end{tabular}
\end{table}


For the sequential image classification tasks (sMNIST, psMNIST, npCIFAR-10), we performed \textbf{100 trials} per model configuration, optimizing for validation accuracy (maximization). For the UCR time series tasks (FordA, FordB, Mackey-Glass), we performed \textbf{100 trials} per configuration, optimizing for validation accuracy (maximization) for classification tasks and NRMSE (minimization) for Mackey-Glass.

All experiments used:
\begin{itemize}
    \item Random seed: 42 (for TPE sampler reproducibility)
    \item Training batch size: 828 (image tasks) / 64 (UCR tasks)
    \item Validation split: 20\% holdout from training data
    \item Readout: Ridge regression (Mackey-Glass) or Logistic regression (classification tasks)
    \item Feature scaling: StandardScaler applied to reservoir activations
\end{itemize}


\subsection{Parameter Budget Matching}
\todo{questa parte dovrebbe essere scritta prima...}
To ensure fair comparison across architectures,
 we employed a parameter matching strategy that adjusts the number of hidden units 
 per layer to achieve approximately 100,000 reservoir parameters.
  The exact number of units depends on:
\begin{itemize}
    \item The number of layers
    \item The input dimensionality (1 for univariate tasks, 96 for npCIFAR-10)
    \item The architecture type (DCR with no recurrence blocks have fewer parameters due to zero recurrence, so more units are used)
\end{itemize}

\begin{comment}
\subsection{Readout Training}

For classification tasks, a logistic regression classifier was trained on the final reservoir state.
 For the Mackey-Glass regression task, Ridge regression was used. In both cases, reservoir activations were standardized using normalization before training the readout.
\end{comment}

\section{Summary of Experiments}

Table~\ref{tab:experiment_summary} summarizes the experimental configurations across all datasets.

\begin{table}[htbp]
\centering
\caption{Summary of experimental configurations by dataset.}
\label{tab:experiment_summary}
\begin{tabular}{@{}lcccccc@{}}
\toprule
\textbf{Dataset} & \textbf{Seq. Length} & \textbf{Input Dim.} & \textbf{Classes} & \textbf{Trials} & \textbf{Metric} & \textbf{Direction} \\
\midrule
sMNIST & 784 & 1 & 10 & 100 & Accuracy & Maximize \\
psMNIST & 784 & 1 & 10 & 100 & Accuracy & Maximize \\
npCIFAR-10 & 1000 & 96 & 10 & 100 & Accuracy & Maximize \\
FordA & 500 & 1 & 2 & 50 & Accuracy & Maximize \\
FordB & 500 & 1 & 2 & 50 & Accuracy & Maximize \\
Mackey-Glass & Variable & 1 & --- & 50 & NRMSE & Minimize \\
\bottomrule
\end{tabular}
\end{table}

\paragraph{Protocol.}
We generate a long sequence, discard $T_{\mathrm{wash}}$ steps, train on $T_{\mathrm{train}}$, validate on $T_{\mathrm{val}}$, and report performance on a disjoint test segment of length $T_{\mathrm{test}}$.

\section{Summary of reported results}
\label{sec:exp:summary}

Table~\ref{tab:exp:summary} summarizes the tasks, outputs, and metrics.

\begin{table}[t]
\centering
\begin{tabular}{lll}
\hline
\textbf{Task} & \textbf{Type} & \textbf{Metric} \\
\hline
Memory Capacity & intrinsic / linear readout & $\mathrm{MC}=\sum_k \mathrm{MC}_k$ \\
Mackey--Glass & regression & NRMSE \\
sMNIST & classification & accuracy (\%) \\
psMNIST & classification & accuracy (\%) \\
npCIFAR-10 & classification / robustness & accuracy (\%), accuracy vs.\ noise \\
\hline
\end{tabular}
\caption{Overview of experimental tasks and evaluation metrics.}
\label{tab:exp:summary}
\end{table}

